% =====================================================================
\chapter{Ángulos, razones trigonométricas y circunferencia unidad}
\label{tema:1}
% =====================================================================

Este tema fija la medida de ángulos en radianes, define las razones trigonométricas primero en el triángulo rectángulo y después, para cualquier número real, mediante la circunferencia unidad, y deduce de las simetrías de la circunferencia las propiedades básicas de signo, paridad y periodicidad. Se sigue el capítulo~7 de~\cite{osat} (§7.1 a §7.4) y el doble enfoque, del triángulo y de la circunferencia unidad, de Stewart, Redlin y Watson~\cite{stewart}.

\section{Ángulos y radianes}

Un ángulo en posición estándar tiene el vértice en el origen y el lado inicial sobre el semieje positivo de abscisas; es positivo si se recorre en sentido antihorario y negativo en sentido horario. Dos ángulos son \emph{coterminales} si tienen el mismo lado final.

\begin{definicion}[Radián]\label{def:radian}
Sea una circunferencia de radio $r>0$ y un ángulo central que abarca un arco de longitud $s$. La medida del ángulo en \emph{radianes} es $\theta=s/r$. Una vuelta completa mide $2\pi$ radianes, de modo que
\[
180^{\circ}=\pi\ \text{rad},\qquad 1^{\circ}=\frac{\pi}{180}\ \text{rad}\approx 0{,}017\,45\ \text{rad},\qquad 1\ \text{rad}=\frac{180^{\circ}}{\pi}\approx 57{,}30^{\circ}.
\]
\end{definicion}

La definición no depende del radio elegido: al multiplicar $r$ por $\lambda$, la semejanza multiplica también $s$ por $\lambda$, y el cociente no cambia. El radián es, por tanto, una razón entre dos longitudes y carece de dimensiones.

\begin{proposicion}[Arco y sector]\label{prop:arco-sector}
En una circunferencia de radio $r$, un ángulo central de $\theta$ radianes ($0\le\theta\le 2\pi$) abarca un arco de longitud $s=r\theta$ y un sector circular de área $A=\tfrac12 r^{2}\theta$.
\end{proposicion}

\addtocounter{definicion}{1}%

\section{Razones trigonométricas en el triángulo rectángulo}

\begin{definicion}[Razones de un ángulo agudo]\label{def:razones-triangulo}
Sea un triángulo rectángulo con un ángulo agudo $\theta$, cateto opuesto $a$, cateto contiguo $b$ e hipotenusa $c$. Se definen
\[
\sin\theta=\frac{a}{c},\qquad \cos\theta=\frac{b}{c},\qquad \tan\theta=\frac{a}{b},
\]
y sus inversas aritméticas $\csc\theta=1/\sin\theta$, $\sec\theta=1/\cos\theta$, $\cot\theta=1/\tan\theta$.
\end{definicion}

\begin{proposicion}\label{prop:razones-bien-definidas}
Las razones de la \cref{def:razones-triangulo} solo dependen de $\theta$, no del triángulo rectángulo elegido.
\end{proposicion}

\begin{proposicion}[Valores notables]\label{prop:notables}
Se cumplen $\sin\frac{\pi}{6}=\cos\frac{\pi}{3}=\frac12$, $\sin\frac{\pi}{3}=\cos\frac{\pi}{6}=\frac{\sqrt3}{2}$ y $\sin\frac{\pi}{4}=\cos\frac{\pi}{4}=\frac{\sqrt2}{2}$.
\end{proposicion}

\section{La circunferencia unidad}

La \cref{def:razones-triangulo} solo tiene sentido para $0<\theta<\pi/2$. Para extenderla a cualquier número real se usa la circunferencia de radio~1 centrada en el origen, $x^{2}+y^{2}=1$ (\cref{fig:circunferencia}).

\begin{definicion}[Seno y coseno de un número real]\label{def:circunferencia}
Para $\theta\in\R$, sea $P(\theta)$ el punto de la circunferencia unidad al que se llega recorriendo sobre ella, desde $(1,0)$, un arco de longitud $\abs{\theta}$, en sentido antihorario si $\theta\ge 0$ y horario si $\theta<0$. Se definen $\cos\theta$ y $\sin\theta$ como las coordenadas de ese punto:
\[
P(\theta)=(\cos\theta,\ \sin\theta).
\]
Cuando $\cos\theta\neq 0$ se define $\tan\theta=\sin\theta/\cos\theta$, y análogamente $\sec$, $\csc$ y $\cot$.
\end{definicion}

\addtocounter{figure}{1}%

\begin{proposicion}\label{prop:coherencia}
Para $0<\theta<\pi/2$, la \cref{def:circunferencia} coincide con la \cref{def:razones-triangulo}.
\end{proposicion}

\begin{teorema}[Propiedades básicas]\label{teo:simetrias}
Para todo $\theta\in\R$ y todo $k\in\Z$:
\begin{enumerate}
  \item (identidad pitagórica) $\sin^{2}\theta+\cos^{2}\theta=1$; en particular, $\abs{\sin\theta}\le 1$ y $\abs{\cos\theta}\le 1$;
  \item (periodicidad) $\sin(\theta+2k\pi)=\sin\theta$ y $\cos(\theta+2k\pi)=\cos\theta$; más aún, $P(\theta)=P(\varphi)$ si y solo si $\theta-\varphi\in2\pi\Z$;
  \item (paridad) $\cos(-\theta)=\cos\theta$ y $\sin(-\theta)=-\sin\theta$;
  \item (suplementario) $\sin(\pi-\theta)=\sin\theta$ y $\cos(\pi-\theta)=-\cos\theta$;
  \item (media vuelta) $\sin(\theta+\pi)=-\sin\theta$ y $\cos(\theta+\pi)=-\cos\theta$;
  \item (complementario) $\sin(\tfrac{\pi}{2}-\theta)=\cos\theta$ y $\cos(\tfrac{\pi}{2}-\theta)=\sin\theta$.
\end{enumerate}
\end{teorema}

\begin{corolario}[Signos por cuadrantes]\label{cor:signos}
El seno es positivo en los cuadrantes I y II; el coseno, en I y IV; la tangente, en I y III.
\end{corolario}

\begin{definicion}[Ángulo de referencia]
El \emph{ángulo de referencia} de $\theta$ es el ángulo agudo $\theta'\in[0,\pi/2]$ que forma el lado final de $\theta$ con el eje $x$. Por el \cref{teo:simetrias}, $\abs{\sin\theta}=\sin\theta'$ y $\abs{\cos\theta}=\cos\theta'$; el signo lo da el \cref{cor:signos}.
\end{definicion}

\addtocounter{definicion}{1}%

\begin{table}[ht]
\centering\small
\caption{Valores notables. Los de los demás cuadrantes se obtienen con el \cref{teo:simetrias}.}\label{tab:notables}
\begin{tabular}{@{}lccccc@{}}
\toprule
$\theta$ & $0$ & $\pi/6$ & $\pi/4$ & $\pi/3$ & $\pi/2$ \\
\midrule
$\sin\theta$ & $0$ & $1/2$ & $\sqrt2/2$ & $\sqrt3/2$ & $1$ \\
$\cos\theta$ & $1$ & $\sqrt3/2$ & $\sqrt2/2$ & $1/2$ & $0$ \\
$\tan\theta$ & $0$ & $\sqrt3/3$ & $1$ & $\sqrt3$ & no definida \\
\bottomrule
\end{tabular}
\end{table}

\addtocounter{definicion}{1}%

\addtocounter{definicion}{1}%

